\section{Frames as an interface}\label{sec:frames}

The model of Section~\ref{sec:model} treats statements as opaque. In practice, most statements in a scientific record are formulas,
and a formula means something only within a frame: the same symbols denote different things for different observers, in different
units, or under different model assumptions. QAM therefore treats a frame as an \emph{interface} with two sides. On the right stands
the formula in ordinary notation, exactly as a physicist or mathematician would write it. On the left stand its roles: the frame, the
language or theory in which it is to be read, and its status. Inside a declared context the left side is implicit and the text looks like
any other paper. Outside that context, or when a formula is carried to another frame, the left side becomes visible.

\subsection{Profiles and contexts}

A \emph{profile} fixes one interpretation within one frame. It consists of a key, a frame, a language, a status and a short description
of the interpretation; once declared it cannot be changed. A \emph{context} selects a profile as the default for a stretch of text.
Formulas registered inside the context are stored with the profile as a premise. Because the profile is hashed together with its
description (Definition~\ref{def:hash}), changing the interpretation changes the hash of every formula registered under it.

The following lines appear in the source of this article. They register the two frames of the example, a profile for each, and the
first formula; the package checks them while the article is typeset.

\begin{lstlisting}
\qamframe{Lab}\qamframe{Boosted}
\qamprofile{lab}{Lab}{L_{\rm EM}}{conditional}{Laboratory fields and their Gram readouts}
\qamprofile{boosted}{Boosted}{L_{\rm EM}}{conditional}{Fields of an observer boosted along z, c=1}
\begin{qamcontext}{lab}
\[ \qamformula{field-pair}{conditional}{F_b=[2e_x\mid b\,e_y],\quad b\in\{+1,-1\}} \]
\end{qamcontext}
\end{lstlisting}

\qamframe{Lab}\qamframe{Boosted}
\qamprofile{lab}{Lab}{L_{\rm EM}}{conditional}{Laboratory fields and their Gram readouts}
\qamprofile{boosted}{Boosted}{L_{\rm EM}}{conditional}{Fields of an observer boosted along z, c=1}
\noindent Inside the context the formula prints as an ordinary display:
\begin{qamcontext}{lab}
\[
\qamformula{field-pair}{conditional}{F_b=[2e_x\mid b\,e_y],\quad b\in\{+1,-1\}}
\]
\[
\qamformula[field-pair]{gram-input}{conditional}{\rho(F_b)=\operatorname{diag}(4/5,1/5)}
\]
\end{qamcontext}
These are the two fields of the example with $t=b$ and $u=0$. Both have the same Gram register, as the second display records.
Outside the context, a reference to the stored formula shows its roles automatically; \verb|\qamphys{gram-input}| prints
\[
\qamphys{gram-input}
\]
so that a reader who meets the formula elsewhere sees in which frame and language it was stated.

\subsection{Explicit conversion}

A formula is carried to another frame only by an explicit step that names the declared map, the source formula, the target profile and
the target formula. The target formula is supplied by the author; QAM does not compute it.

\qammap[\Lambda_z]{lorentz-z}{Lab}{Boosted}
\begin{lstlisting}
\qammap[\Lambda_z]{lorentz-z}{Lab}{Boosted}
\[ \qamphysconvert{boosted-fields}{conditional}{lorentz-z}{field-pair}{boosted}
     {F'_b=\gamma\,[(2-vb)\,e_x\mid(b-2v)\,e_y]} \]
\end{lstlisting}
prints, with both sides of the interface visible,
\[
\qamphysconvert{boosted-fields}{conditional}{lorentz-z}{field-pair}{boosted}
  {F'_b=\gamma\,[(2-vb)\,e_x\mid(b-2v)\,e_y]}
\]
and the Gram registers in the new frame are then registered inside the Boosted context:
\begin{qamcontext}{boosted}
\[
\qamformula[boosted-fields]{gram-output}{conditional}
  {\rho(F'_{+1})\big|_{v=1/4}=\operatorname{diag}(49,4)/53,\qquad \rho(F'_{-1})\big|_{v=1/4}=\operatorname{diag}(9,4)/13}
\]
\end{qamcontext}
The record now states, in a form the checkers can follow, what Section~\ref{sec:example} found by hand: the conversion runs from
\texttt{field-pair}, the fields, and not from \texttt{gram-input}, the register. A conversion from the register would require a map
that does not exist, and nothing in the record claims one.

\subsection{Repairing the map by enriching the description}

The register lost exactly two pieces of information: the overall scale $n=\tr G$ and the orientation of the fields in space. Write
$D=(\rho,n,O)$ with $O\in SO(3)$, and reconstruct the fields as $F=O\,F_0(n\rho)$, where $F_0(G)$ is the upper-triangular
$3\times2$ matrix with $F_0^{\mathsf T}F_0=G$ (a Cholesky factor padded with a zero row). For the two fields above,
$G=\diag(4,1)$, $F_0=[2e_x\mid e_y]$, $n=5$ and $O=\diag(1,b,b)$, so $D_b=(\diag(\frac45,\frac15),5,\diag(1,b,b))$ and indeed
$OF_0=F_b$. The two inputs now differ in their orientation component. The composite ``reconstruct, boost, encode'' is a well-defined
map on such descriptions wherever the boosted field again has rank two, which holds for both fields at $v=\frac14$. In QAM it is
declared as a second map from Lab to Boosted, and the corresponding entries are conversions through it. The normalised output is
unchanged, because a common scale cancels in $\rho$; what has changed is that the record now contains a map that exists.

\subsection{What is checked and what is not}

The package checks that a context is open when a formula is registered, that profiles and formulas are not declared twice, that a
conversion uses a declared map whose source and target match the profiles involved, that two profiles in the same frame are not mixed
without an explicit step, and that status never rises (no target is stronger than its source or its target profile). It does not check
the formulas. They are stored as \TeX{} text, not as expressions of a computer algebra system, so the boost law, units, domains and the
truth of a conversion are not verified. Section~\ref{sec:requests} shows how calculations whose results matter are instead submitted to
an executor and stored with a receipt.
